\textbf{Verificación subsistema 5}
\textbf{Sistema de autorregulación de temperatura
}
Para este sistema se realizaron algunas pruebas, en un primer momento solo se conectó el cinturón para comprobar que este se calentara y observar el comportamiento que presentaba.
Con este sencillo experimento se pudo notar que al calentar, el dispositivo pasa de un color rojo a un color marrón. Posteriormente, se realizaron algunos ensayos colocando tierra en al interior del tambo simulando los desechos vegetales, se procedió a envolver el contenedor con el cinturón y se prosiguió a conectar el cinturón.
Con ayuda de un termómetro se realizaron algunas mediciones de la tierra, previo a encender el cinturón y después de mantenerlo encendido un periodo determinado de tiempo. Al iniciar la prueba la tierra tenía una temperatura de 23°C y después de mantener encendido el dispositivo por un rato se registró una temperatura de 42°C. 

\textbf{Sistema de ventilación}
Durante la fase pruebas del compostador se observaron dos situaciones principalmente , la primera en relación a la salida del vapor generado pro el calentamiento de los desechos, ya que en efecto hay salida de vapor en el filtro puesto que al retirarse dicho componente se puede apreciar cierta condensación del vapor. 
Por otro lado también existe salida de vapor a la entrada del sistema de ventilación, y aunque suene algo contradictorio, esto es algo esperado puesto que cuando el ventilador se encuentra apagado esta zona presenta una menor resistencia al paso del vapor, por solo que este sale sin problema alguno.

Asimismo cuando el ventilador se enciende, se observó una rápida disminución de la temperatura al interior del contenedor, un comportamiento deseable pues permite enfriar el interior de ser necesario y al mismo tiempo introduce oxígeno a la mezcla, es por ello que se propuso incorporarlo al sistema, cabe señalar que en los días fríos pueda disminuir su uso pues es cuando mayores pérdidas de calor se tengan caso contrario a los días calurosos de verano.

\newpage